\chapter{Conclusion} \label{chap:clcsn} The integration of geostatistically associated errors into the base numerical models for the SMaRK methodology was essential for this project, offering significant potential benefits for GIScientists, corporate managers, and policymakers alike. If executed properly, the newly discovered findings of this thesis can be further enhanced with additional in situ measurements at various locations worldwide. The errors observed in the models can be attributed to filling in missing values through interpolation and imputation, the absence of satellite data, and the utilization of a smaller area of interest (AOI). 

Many of the CTM-based transformations applied to available imagery, such as the KE estimates of the DU transformed representation of DAMO$_3$ in Chapter \ref{methods:tomi}, facilitated improved processing of historical databases. Combined with RK, analyses of long-term DAMO$_3$ transport can be efficiently conducted in all available coding languages. While additional monitors can be collected and better interpolated using GEE, along with enhanced project sharing, the preprocessing aspect of this work was designed to establish a framework for ongoing research. The necessary work to accurately fill NA values should be integrated into a larger dissertation aimed at developing a model that effectively encompasses the entirety of the United States. 

The inclusion of different datasets likely mitigated much of the multicollinearity depicted in Figure \ref{fig:allcorrs} which represents the correlation of 28 of the total 47 features discussed in Chapter \ref{chap:tmas}. \begin{figure}[h!]
    \centering
    \includegraphics[width=0.75\linewidth]{ThesisFilePkg//figs/all_28_vars.png}
    \caption{Correlations of All Features Used in All Models}
    \label{fig:allcorrs}
\end{figure} Advanced applications of preprocessing methods noted in this thesis can significantly enhance the continuity of the initial time series for variables such as NDVI, NTL, and CH$_2$O. It is important to note that while interpolation and imputation methods are particularly effective for datasets with only a few days of missing data (e.g., TOM/OMI and ERA5 collections), they may not be suitable for long spans of missing data unless there is careful consideration of historically available data. 

Much of the display and processing was completed in Python and ArcGIS PRO before final importations into GIMP 3.0.2 for final touch-ups and color corrections. Despite the capabilities of ArcGIS Pro and its valuable contributions to the GIS community, the lengthy mapping process involved in this thesis did not justify the additional time required for processing each final raster through a model in ArcGIS for display. Rather, the creation of such a toolbox which automatically corrects satellite imagery given its monitored value should be implemented as a further project and not as the result of an MA thesis.
%%%%%%%%%%%%%%%%%%%%%%%%%%%%%%%%%%%%%%%%%%%%%%%%%%%%%%%%%%%%%%%%%%%%%%%%%%%%%%%%%%%%%%%%%%%%%%%%%%%%%%%%%%%%%%%%%%%%
\section{Feature Evaluation, Importance and Development} \label{clcsn:dev} Most of the satellite data launched in late 2018 were highly correlated with surface ozone. Individual features identified as precursors for surface ozone reactions, based on prior models and simulations (\cite{bowdalo_spectral_2016}; \cite{cai_simulating_2019}; \cite{jacob_effect_1999}), ranked among the top features. The distribution of latitudes, longitudes and elevation were heavily skewed as per Map \ref{fig:basemap} in Chapter \ref{tmas:azbounds}. While a more normal distribution of these features could enhance correlation across various elevations, two primary concerns arise from using latitudes and longitudes as numerical features in geospatial models. Firstly, these coordinates are essentially properties of a geo-atom, such as time, elevation, and site-id; incorporating such properties as features is akin to using a tuning metric within the model. Secondly, the spatial dependency of a monitor's meteorological representation is influenced by a transformation based on its remotely sensed source. 

Numerical representations of spatial data may consequently group these transformations artificially, introducing bias into the models. \begin{figure}[h!]
    \centering
        \centering
        \includegraphics[width=0.75\linewidth]{ThesisFilePkg//figs/hist_model_preds.png}
    \caption{Summary statistics for Historical Features} 
    \label{tab:hist_model}
\end{figure} ML/AI ensembles utilizing geospatial data estimate trends over the defined Area of Interest (AOI) in terms of pixels, each associated with specific spatial data. Figure \ref{tab:hist_model} illustrates that enhancements to CTM-based features through Residual Kriging (RK) in historical data can yield significant results. The limitations imposed by complex trends requiring large volumes of data may be more effectively addressed through the further incorporation and development of these features, particularly in very small AOIs with restricted access to data. RK offers a more stable solution for incorporating geospatial data into ML/AI by treating geospatial data as a separate predictive measure. When the complex trend was appropriately tuned, the estimation of DAMO$_3$ was still more closely aligned with RK integration. 

The remaining histograms not displayed in Chapter \ref{chap:discussion} are presented in Figures \ref{tab:modern_model} and \ref{tab:goat_model}. Spatial information is incorporated into the model without being treated as a feature within the ensemble itself, effectively illustrating spatial dependence and enhancing overall accuracy, as demonstrated in this study. RK is crafted to model spatial drift, which is challenging to capture as a dependent variable in ML/AI models; RK enhancements to the final models address drift while separating predictions from group bias. As per the literature review, finely tuned models such as WRF-CMAQ and MATCH (\cite{dou_wrf-cmaq_2024}; \cite{yu_errors_2018}). \begin{figure}[h!]
    \centering
        \centering
        \includegraphics[width=0.7\linewidth]{ThesisFilePkg//figs/modern_model_preds.png}
    \caption{Summary statistics for Modern Features} 
    \label{tab:modern_model}
\end{figure} Research indicates that incorporating fine resolution land-use data, evaluations of tropospheric/ground level O$_3$ measurements, and employing higher quality modeling techniques can yield similarly accurate results, as evidenced in the ML portions of this thesis (\cite{becker_using_2023}; \cite{chen_comparison_2021}; \cite{hu_estimation_2022}). The integration of geo-statistical methods can more effectively depict the trends for the area of interest (AOI) they are intended to represent. These methods can provide a clearer representation of trends for the AOI.

The methods outlined in this thesis can encompass the boundaries of the United States, utilizing tiled access to census tract resolution DAMO$_3$ values. The transformations and weighting methods utilized in this thesis are indicative of similar transformations employed throughout the history of remote sensing and CTM methods. Many of the best correlating features, as depicted in figure \ref{tab:goat_model} demonstrate that transformation methods exhibit high correlations with their associated counterparts. \begin{figure}[h!]
    \centering
        \centering
        \includegraphics[width=0.75\linewidth]{ThesisFilePkg//figs/goat_model_preds.png}
    \caption{Summary statistics for L.B.D Features} 
    \label{tab:goat_model}
\end{figure} Features such as the DU transformation and KE estimates provide stable depictions of DAMO$_3$, likely due to their ability to remove noise based on thermodynamic concepts. This project found that models yield better results when these features are excluded, as their presence can introduce unintentional bias. When used alongside the CTM feature they represent, they inflate the significance of that feature. This may explain why the models in \ref{tab:hist_model} and \ref{tab:modern_model} exhibit over-performance during training.  

When applied correctly, these types of transformations render older imagery suitable for historical modeling. This thesis enables comprehensive coverage of historical DAMO$_3$ due to imagery generated through remote sensing transformations, such as ERA5, across the United States. It utilizes readily available technology, though patience is necessary with computation times until GPU processing schemes are fully integrated. In general, training the models on various datasets did not consistently improve outcomes or computation times. H.D tended to predict similarly to M.D when using XGRB ensembles with RK. Overall, average RK times would take an additional half to full hour pending the amount of error analyzed and incorporated back into the complex trend.
%%%%%%%%%%%%%%%%%%%%%%%%%%%%%%%%%%%%%%%%%%%%%%%%%%%%%%%%%%%%%%%%%%%%%%%%%%%%%%%%%%%%%%%%%%%%%%%%%%%%%%%%%%%%%%%%%%%%
\subsection{Computation Times} \label{clcsn:comp} Models analyzed in this thesis had vastly different computation times across model tuning, data processing, and RK integration. In general, the random forest (RF) and multi-layered perceptron (MLP) generated promising results; however, each dataset required approximately 2 to 3 hours for adequate training using the random search function. This was not inclusive of RK enhancement, which still took under an hour for these models. Employing a grid search could have allowed each model to identify more precise tuning metrics, albeit at the cost of considerably longer computation times, ranging from 5 to 8 hours depending on the number of features utilized. 

XGRB consistently ran in a matter of seconds, indicating that further tuning and testing are necessary for GPU implementation in RF and MLP methods to achieve reasonable run times. While XGRB was suitable for this project, RF and NNs are more prominent in DAMO$_3$ related literature, as mentioned in Chapter \ref{chap:litreview}. Additionally, proper syntax for the current and future versions of Python was established despite limited knowledge of environment setup and management (e.g., initially using pip and conda simultaneously). Much of the script can be streamlined to enhance computation times on laptops and lower-power machines. This project was developed using Python 3.9, and navigating the various Python versions to Python 3.12.10 was a lesson in itself for the researcher. Updating Python reduced computation times from over 5 hours to just over 2 hours for the final outputs of this project. 

The numerous other libraries that also required updates are likely compatible with Python 3.14 and the upcoming release of 3.15. The desktops used for this research did not incorporate a GPU; one was based on Advanced Micro Devices (AMD) while the architecture of the other was unspecified due to the researcher's lack of system knowledge. Incorporating a GPU requires specialized knowledge and is deferred for the larger AOI. Although this would have significantly reduced computation times as noted in \textcite{zhang2011}, the implementation of AMD's ROCm software was still in development (\cite{ghimire2022}). As software continues to evolve to align with modern technologies, computation times can be drastically reduced through direct incorporation into satellite-based imagery. CUDA is well-suited for implementation in machine learning (ML) and artificial intelligence (AI) ensembles, but it is relatively expensive compared to AMD-based GPU units.
%%%%%%%%%%%%%%%%%%%%%%%%%%%%%%%%%%%%%%%%%%%%%%%%%%%%%%%%%%%%%%%%%%%%%%%%%%%%%%%%%%%%%%%%%%%%%%%%%%%%%%%%%%%%%%%%%%%%%%%%%%%%%%%%
\section{Necessity of Residual Kriging} \label{sec:enhanceIm} The advent of satellite monitoring systems has significantly advanced ozone modeling by supplying remote sensing data essential for large-scale atmospheric studies. The instruments used in this thesis offer invaluable datasets that can be enhanced with monitoring data. Satellite-based datasets are subject to spatial limitations and are typically constrained by their resolution. By themselves, they cannot achieve the spatial and temporal precision needed for accurate, localized exposure assessments. The high-resolution images presented in this thesis call for further investigation into the promising applications of RK in modern numerical models.

By applying residual kriging to baseline satellite data, models can account for local variability and enhance spatial resolution, significantly increasing the dataset's utility for high-resolution historical and current applications. This enhancement is particularly promising for bridging the gap in ozone data quality from 1980 to the present, enabling continuous historical models that track monthly and daily ozone variations across decades. The pursuit of high spatial resolution O$_3$ models may represent a landmark methodology for GIS users. This ambition, driven by the imperative need for readily available and accurate spatial-temporal data, has become attainable due to advances in technology. In particular, the availability of Big Data sources such as TOAR/Google Earth Engine, the integration of spatial uncertainty into machine learning and artificial intelligence representations of satellite data, and the overall improvement of programming language syntax (Python 3.12.x, Java 8.x, R 4.x, SQL Server 2022, etc.) have contributed to this progress. 
%%%%%%%%%%%%%%%%%%%%%%%%%%%%%%%%%%%%%%%%%%%%%%%%%%%%%%%%%%%%%%%%%%%%%%%%%%%%%%%%%%%%%%%%%%%%%%%%%%%%%%%%%%%%%%%%%%%%%%%%%%%%%%%%
\subsection{Enhancement of Satellite Imagery} \label{clcsn:enhanceIm} In tandem, the enhancement of remote sensing monitoring systems through RK underscores the potential for a more comprehensive air pollution monitoring network that combines satellite data with ground-level measurements. This approach addresses a key limitation of remote sensing: the lack of fine-scale spatial detail required for accurate exposure predictions in urban and rural micro-environments. RK can improve the spatial accuracy of satellite data by integrating it with in-situ measurements, thereby creating a more precise depiction of DAMO$_3$ concentrations near densely populated areas. Such high spatial resolution is essential for not only quickly estimating exposures but also for identifying affected areas based on geographic location, socioeconomic status, underlying health vulnerabilities, or a combination of these factors. By refining satellite-based DAMO$_3$ estimates and numerical models, RK can also support more responsive public health interventions, enabling the timely identification and mitigation of high-exposure areas.

Multi-disciplinary GIS modeling techniques such as SMaRK have become increasingly prevalent due to the recent data revolution (\cite{goodchild_mapping_2018}). This revolution is fueled by the exponential growth in data storage and processing capabilities resulting from advancements in Big Data. Consequently, there is a demand for simplified and robust methodologies for the spatial and temporal analysis of both micro- and macro-ecologies. The SMaRK approach has markedly enhanced the predictive accuracy of traditional ML/AI models by spatially interpolating residuals from features at specific points and reintegrating the predicted errors into the model. In the PHOTUC region, satellite imagery leveraged through SMaRK provides transformative capabilities for DAMO$_3$ mapping, CTMs, and exposure studies by incorporating known surface-detected values. Incorporating spatial uncertainty into ML/AI methods is critical for the proper display of geospatial information in GIS programs. 

This integration allows for more detailed and nuanced representations of any independent variable of interest across various temporal and geographic scales, with the appropriate coverage of monitoring stations. These fine-grained enhancements are critical for uncovering the full spectrum of micro- and macro-effects of the many pollutants threatening human health, ecosystems, and air quality. Remote sensing engineers continuously adapt algorithms and address noisy data to reduce errors in their results. The incorporation of monitoring points associated with the processes they depict integrates the concept of Ground Control Points (GCPs) utilized in high-resolution remote sensing through drone technology into modern satallite technologies. GCPs are static points that the drone consistently detects, providing a reference across all images that require tiling. 

High-resolution satellite imagery is stored and shared using data tiling, which requires interpolation and resampling techniques discussed in Chapter \ref{methods:inrtp}. By linking imagery with monitored locations during the resampling process, the concept of uncertain geospatial heterogeneity can be integrated, further reducing image errors. This project demonstrated that applying spatial reasoning in conjunction with machine learning and artificial intelligence ensembles more effectively illustrates geographic patterns and trends identified through the analysis of established observations. With enhanced high-resolution surface DAMO$_3$ representations, numerous communities can be protected from associated risks and their constituents. %%%%%%%%%%%%%%%%%%%%%%%%%%%%%%%%%%%%%%%%%%%%%%%%%%%%%%%%%%%%%%%%%%%%%%%%%%%%%%%%%%%%%%%%%%%%%%%%%%%%%%%%%%%%%%%%%%%%%%%%%%%%%%%%
\subsection{Advancing Surface O$_3$ Mapping with RK Regression} \label{clcsn:advanceRK} The RK method emerged as a prominent approach, providing robust spatial refinement to all baseline models by accounting for residual spatial dependencies. Although some base ensembles outperformed SMaRK-based estimations due to their more complex properties, the overall SMaRK version of these ensembles still demonstrated stronger associations with in-situ values at each location. This approach effectively captured DAMO$_3$ concentrations in various anthropogenic environments with similar meteorological conditions, achieving results equal to or better than current ensemble methods for correcting imagery. 

If the associated instrument can be correlated to a properly gathered monitor of the same type (e.g. Dobson Monitors and TOMS/OMI imagery or high-resolution drone imagery corrected with GCPs and NDVI imagery from Landsat in-between retrievals), then predictive ensembles using remotely sensed data will inherently benefit from the reduced error in imagery. SMaRK’s low error rates and enhanced spatial precision offer novel solutions for correcting high-resolution imagery, with the potential to impact numerous sectors such as public health research, urban chemical transport studies, and local weather mapping, where fine-scale exposure assessments are critical (e.g., \cite{kong_unraveling_2023}, \cite{lo_elevated_2024}, \cite{yin_trend_2019}).

RK can be used to mitigate the Modifiable Areal Unit Problem (MAUP) and offers offers a unique advantage for remotely sensed imagery, providing more accurate and localized predictions based on current and historical technologies. By modeling the relative distance an observation is from its predicted value, the SMaRK approach acts as a mathematical implementation of Tobler’s Law in ML/AI models. Each pixel represents a complex trend based on monitors within a specific field and dictates the amount of similarity between the source and estimated values. By including this in modern imagery, large sets of geo-spatial uncertainty can be accounted for.
%%%%%%%%%%%%%%%%%%%%%%%%%%%%%%%%%%%%%%%%%%%%%%%%%%%%%%%%%%%%%%%%%%%%%%%%%%%%%%%%%%%%%%%%%%%%%%%%%%%%%%%%%%%%%%%%%%%%%%%%%%%%%%%%
\section{Final Comments on the SMaRK Method} \label{clcsn:finalcomments} In general, utilizing the RK method allowed for capturing complex error trends associated with predicted and actual values of surface DAMO$_3$. SMaRK exhibited significant variability and complexity in DAMO$_3$ concentrations due to emissions (\cite{staehle_technical_2023}) while also accounting for subtle influences from meteorological factors (\cite{danyang_ma_effect_2023}; \cite{hou_elucidating_2024}). Although kriging has been recognized for some time (\cite{chiles_fifty_2018}), its frequency of use and applications in geospatial representation within ML/AI ensembles and corrective measures remain limited. Moreover, SMaRK can be integrated into atmospheric monitoring networks established by large institutions, such as the Tropospheric O$_3$ Assessment Report (T.O.A.R.) which already encompasses monitoring and satellite imagery associations (\cite{gaudel_tropospheric_2018}). 

The integration of RK effectively addressed specific limitations inherent in traditional statistical models when managing complex spatial dependencies, thereby enhancing predictions for DAMO$_3$ within the AOI. Statistical models, such as linear and ridge regression, have long been foundational tools for geographic and environmental predictions due to their simplicity and interpretability. However, their effectiveness in modeling DAMO$_3$ is limited by their inability to address non-linear and spatially heterogeneous relationships, which are fundamental characteristics of DAMO$_3$ transport and transformation. Atmospheric chemistry models, such as GEOS-Chem, offer the non-linear predictive power necessary to understand DAMO$_3$ formation processes at large scales, yet they often lack the spatial granularity required for detailed exposure assessments in urban areas. 

Although these models are informative for understanding DAMO$_3$ formation mechanisms, they typically yield results that are too coarse to capture the spatial variations critical to localized public health research. The RK method significantly enhances these models by interpolating spatial residuals from baseline models, thus offering a refined spatial resolution suitable for high-resolution monthly and daily historical modeling. Additionally, improvements to this thesis would facilitate the appropriate tuning of monitor locations in relation to numerous air pollutants mentioned in Chapter \ref{litreview:ctms} and \ref{tmas:S5P}. Further work and the incorporation of this method into spatial programs are necessary to unlock the full potential of GIS in the modern world. 

The establishment of complex trend analysis from GeoAI and Big Data has finally enabled for the proper quantification of surface DAMO$_3$ estimations. Additional applications of RK into statistical frame works could further enhance optical technologies across disciplines that utilize similar observational and correctional frameworks, such as correlating satellite imagery with rovers or stations on Mars. Insights gained from these models can support the development of high-resolution pollutant assessments beyond surface DAMO$_3$ in various fields, providing a robust solution for incorporating geo-statistics into complex ensembles (\cite{cao_deep_2022}; \cite{liu_improve_2018}). Overall, RK represents a fundamental GIScience relationship that can enhance remotely sensed imagery across all domains.
%%%%%%%%%%%%%%%%%%%%%%%%%%%%%%%%%%%%%%%%%%%%%%%%%%%%%%%%%%%%%%%%%%%%%%%%%%%%%%%%%%%%%%%%%%%%%%%%%%%%%%%%%%%%%%%%%%%%%%%%%%%%%%%%
\subsection{Developments and Future Directions} \label{clcsn:developmentanddirections} High spatial resolution DAMO$_3$ models require substantial computational resources and ongoing support for data collection and analysis (\cite{duan_regional_2022}; \cite{rowland_comparative_2024}; \cite{zhou_exposure_2022}). The production of continuous historical models from 1980 to 2000 necessitates significant funding for the acquisition and processing of archival data, as well as the refinement of modeling techniques to ensure accuracy. Establishing daily average models from 2000 to the present will require a robust investment in hardware and software infrastructure to accommodate the high-resolution demands of RK-enhanced datasets. This sustained funding is essential not only for generating accurate DAMO$_3$ models but also for supporting public health research that which impacts policy and environmental regulations.

A commitment to understanding the full scope of DAMO$_3$’s impact on human health extends beyond the scientific community and necessitates collaborative action from public health officials, environmental agencies, and policymakers. The development of high-resolution DAMO$_3$ models could illuminate the toxicology of surface DAMO$_3$ with unprecedented clarity, highlighting the exposure risks associated with both indoor and outdoor environments where monitoring is available. This knowledge is essential for protecting vulnerable populations, particularly those residing in regions frequently affected by DAMO$_3$ pollution events or those working in industries with elevated indoor DAMO$_3$ exposures mentioned in Chapter 2 \ref{litreview:sofat}. DAMO$_3$ modeling that accounts for both indoor and outdoor exposure is crucial for accurately estimating the true burden of DAMO$_3$ related diseases, as indoor exposure can vary significantly depending on building ventilation, material interactions, and proximity to outdoor pollution sources (\cite{jadoon_health_2022}; \cite{barzeghar_long-term_2020}; \cite{hwang_indoor_2020}). A comprehensive understanding of DAMO$_3$ exposure patterns could facilitate policies aimed at reducing indoor DAMO$_3$ levels in workplaces, schools, and homes, thus further protecting public health.

Accurate DAMO$_3$ models have the potential to reduce mortality rates associated with air pollution by informing more targeted and effective interventions (e.g., in \cite{li_spatiotemporal_2024} and \cite{rovira_air_2020}). By providing a clearer picture of where and when high DAMO$_3$ concentrations occur, these models empower public health officials to issue more precise advisories and allocate resources effectively to areas in need. This can lead to reduced hospital admissions, fewer cases of respiratory distress, and ultimately a decrease in the mortality rate associated with DAMO$_3$-related health issues. The insights gained from high-resolution DAMO$_3$ models can also serve as a foundation for preventative measures, such as urban planning initiatives that aim to reduce vehicle emissions in densely populated areas, minimizing the risk of acute DAMO$_3$ exposure. Such measures not only enhance air quality but also contribute to sustainable urban environments, promoting public health and well-being in the long term.
%%%%%%%%%%%%%%%%%%%%%%%%%%%%%%%%%%%%%%%%%%%%%%%%%%%%%%%%%%%%%%%%%%%%%%%%%%%%%%%%%%%%%%%%%%%%%%%%%%%%%%%%%%%%%%%%%%%%
\subsection{Omission of Certain Methods and Data Assumptions} \label{clcsn:comp} Given the scale of the final predictive value, Box-Cox or Yeo-Johnson transformations were not utilized. During the initial drafting of this thesis, the use of normalized values resulted in minimal improvements in the overall accuracy of each model and inadvertently introduced confusion by maintaining similar value ranges across models. Scaling the dataset yielded the same deviations as unscaled data without requiring an additional step in an otherwise arduous process. While these transformations are typically used for small datasets, Box-Cox transformation requires all positive values within the feature set; thus, updates were made using the Yeo-Johnson method, which accommodates both positive and negative observations (\cite{riani2023}).

Development of larger tiles, the inclusion of additional monitoring systems, and the depiction of historical trends are reserved for future work by the researcher. Regarding census data and statistics, population data for the years 2024 and onward were not used, as future updates like these are reserved for subsequent work. Daily population data was not constructed as the frame work for this is incorporated into the script for DAMO$_3$ by way of "datetime" objects and are also meant for future work. This project also required the assumption that monitors were not moved, removing such monitors from the AOI. If a monitor is moved, it loses the ability to provide insights into potential future trends unless its original location and measurements are accurately reintegrated before initialization within an ensemble. 
%%%%%%%%%%%%%%%%%%%%%%%%%%%%%%%%%%%%%%%%%%%%%%%%%%%%%%%%%%%%%%%%%%%%%%%%%%%%%%%%%%%%%%%%%%%%%%%%%%%%%%%%%%%%%%%%%%%%%%%%%%%%%%%%
\section{Closing Remarks} \label{clcsn:closing} Atmospheric chemistry plays a fundamental role in the history of life on Earth (\cite{barzeghar_long-term_2020}; \cite{chapleski_heterogeneous_2016}; \cite{flynn_spatial_2021}; \cite{tanimoto_diagnosing_2008}). Abiogenetic processes occurring nearly 4 billion years ago were linked to the emergence of oxygen during the Great Oxidation Event (GOE), which enabled the formation of the O$_3$ layer and promoted photosynthetic processes (\cite{schell_reformulated_2002}; \cite{lyons_rise_2014}). As O$_3$ and organisms grew in abundance and complexity, atmospheric composition began to stabilize, ultimately giving rise to the Earth as we know it today. For eons, the specific composition of chemicals on Earth has facilitated the evolution of life by maintaining hypoxia tolerances through variations in energy levels reaching the surface (\cite{sperling_breathless_2022}).  

The safety and well-being of past, present, and future generations is fundamental to the work of many in geography. Applications of spatial analytics are utilized across various disciplines to refine and theorize the complexities of human health and cognition, thereby enhancing our understanding of environmental impacts on daily life (\cite{jerrett2009}; \cite{greenough_beyond_2019}; \cite{nawaz_source_2023}). Decades of integrating geospatial computer science with statistical modeling have resulted in the emergence of unique, multidisciplinary specializations that explore these topics. As computer science continues to improve in availability, quality, and, most importantly, accessibility, geographers increasingly incorporate complex spatial thought into “sibling sciences,” fostering innovative solutions to both known and unknown problems (\cite{goodchild_sharing_2007}). 

To reduce DAMO$_3$-related health risks, support sustainable urban environments, and restore the Earth’s natural DAMO$_3$ cycles, the SMaRK method facilitates a more accurate representation of DAMO$_3$ concentrations across diverse and homogeneous settings. Enhancing current modeling efforts for general air pollution has become increasingly necessary due to ongoing concerns regarding known and emerging health outcomes (\cite{im_present-day_2023}; \cite{sofen_how_2015}; \cite{yan_impacts_2023}). Developing theories on health associations stemming from surface DAMO$_3$ exposure among pediatric and elderly populations necessitates more detailed DAMO$_3$ models to efficiently assign accurate exposures that can support air quality, public health, and urban ecology studies (\cite{deng_effects_2023}; \cite{li_double_2023}; \cite{shin_air_2018}). Given that data for these studies often covers limited areas compared to entire nations or states, there is a pressing need for methodologies that integrate machine learning with geographic information. This integration can fully leverage the benefits of both techniques for high-resolution geographic modeling. 

%%%%%%%%%%%%%%%%%%%%%%%%%%%%%%%%%%%%%%%%%%%%%%%%%%%%%%%%%%%%%%%%%%%%%%%%%%%%%%%%%%%%%%%%%%%%%%%%%%%%%%%%%%%%%%%%%%%%%%%%%%%%%%%%
\subsection{Accessing Code and Data} \label{clcsn:access} All raster data is freely available on GitHub via the following link: \url{https://github.com/Ryan-Erickson911/surface_ozone}. Scripts for downloading the necessary data are written in Python. While the complete project is privately available on GitHub, the public version includes access to scripts and sample data which allow for access to the materials mentioned in Chapter \ref{chap:tmas}. The full scope of this project remains in its early stages. Navigating the available code may currently pose challenges; however, ongoing developments and improvements are being implemented daily to create a more user-friendly Python package aimed at facilitating the development of high-spatial-resolution imagery.